% The γ slide's figure (user, 2026-10-29, mock-up option C): SPY and the market state ξ over the real window,
% Jan 2, 2025 – Sep 4, 2026, from talk/data/market_state.csv and market_state_crossings.csv (scripts/30_talk_data.jl,
% which checks that every series stays inside its strip's y range and that each crossing is a change in the EMAs'
% order). Input once in the preamble, after vnslides (for vnplotpanel, vnplotgrid and vnband) and
% stylized_facts.tex (for the vnsf look).
%
% \vnmarketstate{<width>}{<SPY height>}{<ξ height>} draws two strips, one above the other, each with a plot area
% <width> wide:
%   - SPY's close in faded gold (gold, as SPY is on every slide), under its 21-day EMA (white) and 63-day EMA (blue),
%     all three ÷ SPY's first close; a key at the bottom middle, August 2025 to February 2026, where no day is shaded
%     and every line is above 1.03;
%   - ξ = −G(EMA21/EMA63 − 1), with its positive part filled, labelled at −1, 0 and 1; its zero line is named
%     EMA21 = EMA63 at the right end, where ξ stays below −0.2.
% The two EMAs are about half a millimetre apart on the strip (user: the crossovers are too hard to see), so each
% crossing gets an open ring on the EMAs, the lines visible through it, and a dotted leader down to a dot where ξ
% crosses zero. Both strips shade the days with ξ > 0 (the data's `falling` column), so the sell-offs line up. x is
% the calendar date, ticked on the first of each quarter and titled under the ξ strip; the first date starts at the
% left edge, clear of the −1 beside it.

\usepgfplotslibrary{dateplot}

% the date axis and the shaded days, shared by both strips
\pgfplotsset{
  vnms/.style={
    vnsf, width=#1, xmin=2025-01-01, xmax=2026-09-04, date ZERO=2025-01-01, date coordinates in=x,
    xtick={2025-04-01,2025-07-01,2025-10-01,2026-01-01,2026-04-01,2026-07-01},
    xticklabel style={/pgf/number format/assume math mode=false},
  },
}

% \vnmsshade{<low>}{<high>}: the days with ξ > 0, filled from <low> to <high> on the strip's y axis
\newcommand{\vnmsshade}[2]{%
  \addplot[name path=vnmstop, draw=none, const plot, forget plot]
    table[col sep=comma, x=date, y expr={#1 + (#2 - #1)*\thisrow{falling}}] {data/market_state.csv};
  \addplot[name path=vnmsbot, draw=none, forget plot] coordinates {(2025-01-01,#1) (2026-09-04,#1)};
  \addplot[fill=vnband, fill opacity=0.16, draw=none, forget plot] fill between[of=vnmstop and vnmsbot];
}

\newcommand{\vnmarketstate}[3]{%
  \pgfplotstableread[col sep=comma]{data/market_state_crossings.csv}\vnmsCross
  \pgfplotstablegetrowsof{\vnmsCross}%
  \pgfmathtruncatemacro\vnmsLastRow{\pgfplotsretval - 1}%
  \begin{tikzpicture}
    % SPY and its EMAs, scaled to 1 on the first day; no tick labels at the edges, so the corners stay clear
    \begin{axis}[vnms=#1, name=vnmsspy, height=#2,
      xticklabels={},
      ymin=0.8, ymax=1.4, ytick={0.9,1.1,1.3},
      ylabel={SPY},
      legend columns=3,
      legend style={at={(rel axis cs:0.53,0.03)}, anchor=south, fill=vnplotpanel, draw=none, inner sep=1pt,
        font=\fontsize{7}{8.4}\selectfont, text=vndottick, /tikz/every even column/.append style={column sep=7pt}},
      % ## is the line style pgfplots passes in, not the macro's first argument
      legend image code/.code={\draw[##1] (0cm,0cm) -- (0.42cm,0cm);},
    ]
      \vnmsshade{0.8}{1.4}
      % SPY faded, so the EMAs read
      \addplot[draw=vnfigaccent, draw opacity=0.55, line width=0.6pt, line join=round, mark=none]
        table[col sep=comma, x=date, y=spy] {data/market_state.csv};
      \addplot[draw=vnslideink, line width=1.0pt, line join=round, mark=none]
        table[col sep=comma, x=date, y=ema_short] {data/market_state.csv};
      \addplot[draw=vnband, line width=1.0pt, line join=round, mark=none]
        table[col sep=comma, x=date, y=ema_long] {data/market_state.csv};
      \legend{SPY, EMA$_{21}$, EMA$_{63}$}
      % each crossing on the EMAs, by row (from 0) of the crossings file; pgfplots runs drawing commands at
      % \end{axis}, so \edef fixes the date and level now
      \foreach \vnmsR in {0,...,\vnmsLastRow} {
        \pgfplotstablegetelem{\vnmsR}{date}\of\vnmsCross
        \let\vnmsDate\pgfplotsretval
        \pgfplotstablegetelem{\vnmsR}{level}\of\vnmsCross
        \let\vnmsLevel\pgfplotsretval
        \edef\vnmsCmd{\noexpand\coordinate (vnmsup\vnmsR) at (axis cs:\vnmsDate,\vnmsLevel);}%
        \vnmsCmd
      }
    \end{axis}
    % the market state
    \begin{axis}[vnms=#1, name=vnmsxi, height=#3,
      at={(vnmsspy.south west)}, anchor=north west, yshift=-0.22cm,
      xticklabels={Apr,Jul,Oct,Jan 2026,Apr,Jul},
      extra x ticks={2025-01-01}, extra x tick labels={Jan 2025},
      extra x tick style={grid=none, xticklabel style={anchor=north west, xshift=-1pt}},
      ymin=-1, ymax=1, ytick={-1,0,1}, yticklabels={−1,0,1},
      ylabel={$\xi$},
      xlabel={date},
    ]
      \vnmsshade{-1}{1}
      \addplot[name path=vnmspos, draw=none, forget plot]
        table[col sep=comma, x=date, y expr={max(\thisrow{xi}, 0)}] {data/market_state.csv};
      \addplot[name path=vnmszero, draw=none, forget plot] coordinates {(2025-01-01,0) (2026-09-04,0)};
      \addplot[fill=vnband, fill opacity=0.45, draw=none, forget plot] fill between[of=vnmspos and vnmszero];
      \addplot[draw=vnslidelink, line width=0.8pt, line join=round, mark=none]
        table[col sep=comma, x=date, y=xi] {data/market_state.csv};
      % the zero line, named: ξ is zero exactly where the two EMAs cross
      \node[anchor=south east, inner sep=1.5pt, font=\fontsize{7}{8.4}\selectfont, text=vndottick]
        at (rel axis cs:1,0.5) {EMA$_{21}$ = EMA$_{63}$};
      % each crossing on ξ's zero line
      \foreach \vnmsR in {0,...,\vnmsLastRow} {
        \pgfplotstablegetelem{\vnmsR}{date}\of\vnmsCross
        \let\vnmsDate\pgfplotsretval
        \edef\vnmsCmd{\noexpand\coordinate (vnmsdown\vnmsR) at (axis cs:\vnmsDate,0);}%
        \vnmsCmd
      }
    \end{axis}
    % each crossing: an open ring on the EMAs, a dotted leader from the ring down to ξ's zero line, a dot there
    \foreach \vnmsR in {0,...,\vnmsLastRow} {
      \draw[vnslideink, line width=0.5pt, dash pattern=on 0.8pt off 1.4pt]
        ([yshift=-3.2pt]vnmsup\vnmsR) -- (vnmsdown\vnmsR);
      \node[circle, draw=vnslideink, line width=0.9pt, minimum size=6.4pt, inner sep=0pt] at (vnmsup\vnmsR) {};
      \fill[vnslideink] (vnmsdown\vnmsR) circle (1.3pt);
    }
  \end{tikzpicture}%
}
